% !TEX root = ../main.tex
\chapter{The blocks at a glance}\label{ch:blocks}

Every analysis the app can do is a \emph{block}: one step that takes something in, does
one job, and passes something on. You build an analysis by putting blocks in a row, a
\emph{chain}, on the Chain page (\apppage{analyse.chain}). This chapter gives each of the
\ubnumblocks{} blocks a short card in everyday words. Skim it to see what exists; come back
when you need one.

\begin{plainterms}[Blocks are like kitchen tools]
A chain is a recipe and each block is one tool in it: a sieve, a knife, a scale. Each tool
expects a certain kind of thing (you cannot weigh soup with a sieve) and hands on a certain
kind of thing. The app checks the hand-offs for you: if a block cannot take what the step
before it gives, the chain marks that junction invalid and says which kind each side expects,
in the form ``Seed matches needs Signal · previous emits SpanSet''.
\end{plainterms}

\section*{How to read a card}
Each card has the same parts.
\begin{itemize}
\item \emph{The title} is the block's name as the app shows it.
\item \emph{The coloured tag} is the tab of the Insert-a-stage window it sits under (the six
  sections of this chapter follow those tabs).
\item \emph{Takes \ldots, gives \ldots} names the kind of thing the block takes and the kind
  it hands on: a Signal, a SpanSet, a WindowSet, Scores, an Encoding, a Grouping or a Model.
  \Cref{ch:core-ideas} explains each kind.
\item \emph{What it does, When to use it, What goes in and what comes out, Settings that
  matter, Watch out for}: five short answers in plain words.
\item \emph{Full details} links to the block's entry in \cref{app:blocks}: the algorithm,
  every setting with its default, how it fails, how it was checked, and where its code lives.
\end{itemize}

\begin{underthehood}
A card's title, tag, input and output types and appendix link are read from the block's
\code{AdapterSpec} in \code{Adapters/} by \code{scripts/manual_tables.py}, as is the whole
settings table of its appendix entry; the bridge checks the same input and output types
when it validates a chain (\code{webui/server/chain.py}). A block that is registered but
cannot run on this machine says so on its card.
\end{underthehood}

% One paragraph under each group's heading.
\ubcategoryintro{preprocess}{These blocks change the signal before anything looks for
patterns in it: they remove slow drift, keep or drop bands of frequencies, or flip the sign.
\blockref{preprocessing.decompose} splits a trace into its slow baseline, its fast noise and
the signal left between them, and lets you pass any one of the three on.
Two work on windows rather than whole signals: \blockref{preprocessing.trace_shape} reduces
each window to its shape, and \blockref{preprocessing.window_pool} starts a chain from window
sets you saved earlier instead of from a span.}

\ubcategoryintro{encode}{These blocks re-describe the signal in another form: a string of
letters, an image, a spectrum, a set of wavelet coefficients, a fingerprint of numbers. An
encoding does not mark anything by itself; later blocks search it, compare it or learn from
it.}

\ubcategoryintro{detect}{These blocks find where something happens. Most return the
stretches of time where it happens (a span set). Four return a score at every moment
instead (\blockref{detection.matrix_profile}, \blockref{detection.fluss},
\blockref{detection.wavelet_summation} and \blockref{detection.deconvolve}), and four more
turn such scores into stretches (\blockref{detection.threshold},
\blockref{detection.mp_motifs}, \blockref{detection.summation_threshold} and
\blockref{detection.impulse_events}).}

\ubcategoryintro{cluster}{These blocks cut a recording into windows, measure them, and group
the ones that look alike. They make the raw material of the Library and of every model:
windows, the measures on them, and groupings of them, whether drawn by a person or found by
the machine.}

\ubcategoryintro{model}{These blocks learn from labelled windows and then score new ones.}

\ubcategoryintro{control}{These blocks measure what was found and supply the controls to
judge it against: the shape of each event, the intervals between events and the trains
they form, and surrogate data that keeps some properties of a recording while destroying
others, so a pattern can be tested against chance. The one-block version of the Dehshibi
spike detector, \blockref{detection.dehshibi_spikes}, also sits under this tab.}

\ubcatalogue{card}
